\chapter{Introduction}
\label{chap:introduction}

% ===================== GUIDELINE AUDIT 2026-07-25 -- INTRODUCTION =====================
% Checked against the SCSS Report Guidelines and the Dissertation Marking Sheet (weights:
% problem/motivation 10%, background/LR 15%, technical execution 50%, testing/evaluation/
% conclusions 15%, presentation 10%). Full audit: docs/TCD_Dissertation/GUIDELINE_COMPLIANCE.md
%
% This chapter PASSES the guideline's structural requirements: it states the project and its
% purpose, and Section~\ref{sec:dissertation-structure} is the required "reader's guide to the
% report, chapter by chapter". Four fixes remain, listed below and at their sites.
%
% (1) [presentation 10%] NO FIGURE IN THIS CHAPTER. The 70+ marking band asks for "excellent use
%     of figures, tables, and diagrams where appropriate", and the Introduction and Conclusion are
%     currently the only two chapters with neither. fig:architecture is ALREADY DRAFTED below
%     (Section~\ref{sec:research-experiments}) and the supervisor is recorded as liking it.
%     ACTION: uncomment the \begin{figure}...\end{figure} block for fig:architecture AND the
%     [OPENER] sentence above it that references it. This is the single highest-value edit in the
%     chapter. Pick ONE of the three drafted figures for the Motivation (fig:alice-on-device or
%     fig:on-device-overview) if a second figure is wanted -- not both.
%
% (2) [background/LR 15%] CITATION MISATTRIBUTION, still live at Section~\ref{sec:scope-and-assumptions}.
%     The [AUDIT-FIX 2026-07-24] block there gives the corrected wording; it has not been applied.
%     The same wrong claim is live in three further places -- background.tex (the paragraph
%     beginning "Reinforcement learning is set aside on purpose"), state-of-the-art.tex (the
%     paragraph beginning "Models below a billion parameter class"), and conclusion.tex
%     (Section~\ref{subsec:conclusion-future-grpo}). Apply the fix at ALL FOUR sites. A reviewer who
%     opens Rainone et al. finds the claim reversed; under "systematic and complete reference to
%     sources used" this is the most damaging single item in the dissertation.
%
% (3) [presentation 10%] THIRD PERSON. The guideline states "Use third person as much as possible".
%     Two live first-person slips in this chapter, both fixed inline below.
%
% (4) [presentation 10%] CONTRACTIONS. The guideline states "Avoid use of contractions". Two live
%     instances in this chapter, both fixed inline below.
% =====================================================================================

Society is built on trust. People trust that the food they buy is safe to eat, that the cars they drive will not fail without warning, and that those around them will mostly act in good faith. As artificial-intelligence (AI) systems become woven into daily life, they too must earn that trust: AI already helps people with curated news, health advice, financial management and even relationships. The common assumption is that assistance this trustworthy and this personalised requires cloud-scale computation, yet that computation is now shifting onto the user's own device. This dissertation tests whether a small model running entirely on that device can meet the same safety and personalisation standards as its larger, cloud-based counterparts. The study finds a measurable trade-off between constitutional compliance and reasoning capability at this scale, where on-device inference is feasible but behaviour is brittle in ways larger models are not, as Rainone et al.\ observed \cite{rainone2025replacingthinkingtoolusage}. The goal is not a production-ready model for the general public but a reusable method for measuring and improving that compliance on small models. This chapter sets out the motivation, problem statement, research question, hypotheses, experiments, scope, contributions and structure of the dissertation.

\section{Motivation}
\label{sec:motivation}

A shift is under way in how people seek knowledge online, and the search engine is losing its monopoly on the opening question. A 2026 Bain survey finds that only 42\% of millennials and Gen~Z now reach for a search engine first, against 76\% of baby boomers \cite{bain2026snapchart}, and Google's own conversational ``AI Mode'', a reworking of its twenty-seven-year-old search product, passed a billion monthly users within a year of launch \cite{reid2026aisearch}. What replaces the keyword query is a different act: the search engine remains useful for a quick lookup, but the reason people increasingly give for turning to an assistant is to make sense of complex matters through personalised, step-by-step guidance and brainstorming \cite{bain2026snapchart}, and the topics worked through this way now include the health, financial and personal questions once reserved for a professional or a trusted confidant. A personal assistant lets a user ask, clarify and ideate without fear of being judged, and because it is asked things a person might hesitate to put to another human or type into a search box, the record it accumulates is unusually intimate; that is what makes the data it holds sensitive. If a search query captures the \emph{what}, \emph{when} and \emph{where} of an event in a user's life, the conversation captures the \emph{what}, \emph{why} and \emph{how} of the thinking behind it, and it is this record of the user's thought process that constitutes the special-category data GDPR Article~9 singles out for protection \cite{gdpr2016} and that the EU AI Act places in its high-risk tier \cite{europeanparliament2024regulationeu20241689}.

The difficulty in delivering a personalised assistant while keeping the data under the user's control is architectural rather than contractual. A cloud-hosted model requires that data to leave the device on every turn, and the usual safeguards, zero-retention policies, encrypted transit and data-processing agreements, govern what a provider stores, not what a model has already absorbed and can be induced to repeat through targeted queries, as Carlini et al.\ demonstrate \cite{carlini2021extractingtrainingdatalarge}. Federated learning removes the transmission but not the residue, since the trained weights still carry each user's influence, and erasing it, the Article~17 right to be forgotten, remains an open problem, as Nguyen et al.\ and Staufer show \cite{nguyen2026computationcommunicationefficientfederated, staufer2025llmsforgetquantifyingpersonal}. Running the model entirely on the user's own device dissolves both issues at once: nothing is transmitted, and no shared model ever encodes the user, the property the European Data Protection Supervisor identifies as the central data-protection advantage of on-device processing \cite{edps2025ondeviceai}. The largest mobile platforms have reached the same conclusion: in 2026 Apple released a three-billion-parameter on-device Foundation Model, paired with Private Cloud Compute for heavier requests \cite{apple2026afm3}, and Google shipped phone-class Gemma~4 variants of two to four billion parameters capable of running entirely offline \cite{gemma4}. To iterate faster and to reach the low-end and older handsets those flagships leave behind, this dissertation works with a smaller 0.6-billion-parameter model and asks whether a framework can still make it trustworthy. What that choice costs is the rest of the story.

% [DRAFT 2026-07-18 -- Motivation figure built from the EDPS "A day in Alice's life" scenario (TechSonar
% 2025, on-device AI, \cite{edps2025ondeviceai}). It dramatises the privacy argument: today a person's most
% sensitive data leaves the device and can be shared without consent and breached; this work keeps it on the
% device as a single copy. This is a NARRATIVE alternative to the system-overview figure fig:on-device-overview
% drafted before \section{Definitions}; pick ONE for the Motivation (this reads as motivation, the other as an
% architecture overview). To place: uncomment and add "(Figure~\ref{fig:alice-on-device})" to a sentence such as
% the one above about data leaving the device. Check the layout in Overleaf and nudge coordinates as needed.
% \begin{figure}[t]
%   \centering
%   \begin{tikzpicture}[>=Stealth, font=\scriptsize,
%       box/.style={draw, rounded corners, fill=black!5, align=center, text width=3.0cm, minimum height=1.5cm, inner sep=6pt},
%       cloud/.style={draw, rounded corners, fill=black!12, align=center, text width=2.7cm, minimum height=1.5cm, inner sep=6pt},
%       bad/.style={draw, rounded corners, fill=red!8, align=center, text width=2.9cm, minimum height=1.5cm, inner sep=6pt},
%       note/.style={font=\scriptsize\itshape, align=center}]
%     % --- Row 1: the status quo (Alice's story) ---
%     \node[anchor=west, font=\bfseries] at (-0.4,4.2) {Today: personal data leaves the device (EDPS ``A day in Alice's life'')};
%     \node[box]   (dev1)   at (1.9,2.6)  {Alice's devices\\[2pt] \scriptsize sleep \& vitals, voice, location, diet, calendar};
%     \node[cloud] (cloud1) at (7.0,2.6)  {Vendors / cloud / data brokers};
%     \node[bad]   (breach) at (11.8,2.6) {Breach: profile shared \& leaked online};
%     \draw[->, thick, red!70!black] (dev1)   -- (cloud1);
%     \draw[->, thick, red!70!black] (cloud1) -- (breach);
%     \node[note, text width=2.4cm] at (4.55,3.45) {sent off-device, often without consent};
%     % --- Row 2: this dissertation (on-device) ---
%     \node[anchor=west, font=\bfseries] at (-0.4,0.2) {This dissertation: the same data stays on the device};
%     \draw[very thick, rounded corners] (0.2,-3.8) rectangle (8.8,-1.7);
%     \node[anchor=north west, font=\scriptsize\itshape] at (0.35,-1.78) {user's device};
%     \node[box] (dev2) at (2.1,-2.9) {Alice's data\\[2pt] \scriptsize one copy, on device};
%     \node[box] (asst) at (6.7,-2.9) {On-device assistant\\[2pt] \scriptsize Thinker $+$ Executor};
%     \draw[<->, thick] (dev2) -- (asst);
%     \node[cloud, text width=2.5cm] (cloud2) at (11.7,-2.9) {Cloud / vendors};
%     \draw[->, dashed, thick] (8.8,-2.9) -- (cloud2.west);
%     \node[red!75!black] at (9.6,-2.9) {\large\bfseries $\times$};
%     \node[note, text width=2.2cm] at (9.75,-2.05) {nothing transmitted};
%   \end{tikzpicture}
%   \caption{Motivation, after the EDPS ``A day in Alice's life'' scenario \cite{edps2025ondeviceai}: in today's cloud model a person's most sensitive data --- sleep and vitals, voice, location, diet --- is sent off the device, can be shared without consent, and is exposed in a breach; the architecture of this dissertation keeps that data on the device as a single copy, with nothing transmitted to the cloud at inference.}
%   \label{fig:alice-on-device}
% \end{figure}

That cost is more than raw accuracy. When a 0.6B model is fine-tuned to follow constitutional principles, safety behaviour, tool calling and step-by-step thinking together, its limited parameters cannot hold all of these at once, and training for safety degrades the very reasoning it was meant to support. Whether the cause is the training data or the scale itself is not yet settled, the brittleness of models at this scale having been tied by Rainone et al.\ to the interaction of small scale with free-form chain-of-thought reasoning, and shown by them to be partly recoverable through a more structured reasoning format \cite{rainone2025replacingthinkingtoolusage}; but either way it sets a measurable ceiling on a single small model and raises the question the rest of the work turns on: can the lost capability be recovered without making the model larger? Because the framework is developed on, and tested against, a single small model while being designed to transfer independently of it, swapping in a different model in future remains supported.
%  ACTION: CONSIDER (judgement call, not mechanical): in the Motivation (the paragraph beginning "The difficulty in achieving a personalised AI...") the phrase "it dissolves both issues at once" (transmission and shared-model residue) is a strong but architecturally defensible claim about on-device processing; if a plainer register is preferred, "it removes both at once" reads less absolute. (Borderline, likely fine to keep.)
%  (Context: one genuine over-claim in this paragraph; the rest of the Introduction reads as suitably measured, so this is the only tone flag raised here.)

% [DRAFT 2026-07-17 -- introduction figure ("Figure 1.1"). The Motivation argues the whole thesis in one
% picture: an assistant where nothing leaves the phone. A conceptual overview here orients the reader and
% grounds the three pillars of the research question (privacy-preserving, personalised, trustworthy) before
% the definitions. It is deliberately higher-level than the detailed turn-flow figure in the Experiments
% chapter (Figure~\ref{fig:thinker-executor}) -- that one shows the mechanism; this one shows the device
% boundary. Uses only tikz libraries already loaded (positioning, arrows.meta, calc, fit, shapes.geometric).
% To place: uncomment and append "(Figure~\ref{fig:on-device-overview})" to the sentence "If the model run
% entirely on the user's own device, it dissolves both issues at once" in the Motivation above. Coordinates
% are a first pass -- check the layout in Overleaf and nudge as needed.
% \begin{figure}[t]
%   \centering
%   \begin{tikzpicture}[>=Stealth, font=\scriptsize,
%       comp/.style={draw, rounded corners, fill=black!5, align=center, minimum height=0.95cm, inner sep=3pt},
%       store/.style={draw, rounded corners, fill=black!9, align=center, minimum height=0.95cm, inner sep=3pt},
%       note/.style={font=\scriptsize\itshape, align=center}]
%     % device boundary
%     \draw[very thick, rounded corners] (-0.3,0.2) rectangle (12.4,5.0);
%     \node[anchor=north west, font=\footnotesize\bfseries] at (-0.2,4.9) {User's device (offline-capable)};
%     % cloud outside the boundary, faded
%     \node[comp, fill=black!14, minimum width=2.8cm] (cloud) at (9.6,6.1) {Cloud model / server};
%     \draw[->, dashed, thick] (6.2,5.0) -- (cloud.south);
%     \node[red!75!black, font=\normalsize\bfseries] at (7.5,5.55) {$\times$};
%     \node[note] at (4.1,5.55) {no user data ever leaves the device};
%     % components inside
%     \node[comp,  minimum width=2.0cm] (user) at (1.5,3.3) {User\\ conversation};
%     \node[comp,  minimum width=3.0cm] (thk)  at (5.8,4.0) {Thinker 0.6B\\ constitution in weights};
%     \node[comp,  minimum width=3.0cm] (exe)  at (5.8,2.4) {Executor 0.6B\\ tool execution};
%     \node[store, minimum width=2.7cm] (mem)  at (10.3,4.0) {User-model graph\\ (5W+H beliefs)};
%     \node[store, minimum width=2.7cm] (scr)  at (10.3,2.4) {Scratchpad\\ (turn context)};
%     \node[comp,  minimum width=4.4cm] (tools) at (5.8,1.0) {Local tools: Python, date-time, memory};
%     % flows
%     \draw[->]  (user) -- (thk);
%     \draw[->]  (thk)  -- node[note, right, xshift=1pt]{plan} (exe);
%     \draw[<->] (thk)  -- (mem);
%     \draw[<->] (exe)  -- (scr);
%     \draw[->]  (exe)  -- (tools);
%     \draw[->]  (exe.west) to[bend left=18] node[note, below]{response} (user.south);
%   \end{tikzpicture}
%   \caption{Overview of the on-device system this dissertation builds. A Thinker and an Executor (each a
%     0.6B model), the user-model graph, the scratchpad and the local tools all run on the user's own device,
%     so no part of the user's conversation is transmitted to the cloud. The figure grounds the three pillars
%     of the research question: privacy-preserving (the device boundary), personalised (the user-model graph)
%     and trustworthy (the constitution encoded in the Thinker's weights).}
%   \label{fig:on-device-overview}
% \end{figure}
\section{Definitions}
\label{sec:definitions}

Before covering the research question and hypotheses, there are several terms one should know in general in a sense specific to this dissertation. They are defined as follows.

\begin{description}
  \item[Language model.] By definition, a language model is a computational model which predicts the next word in a given sequence of text. Language models vary in different sizes, the models in general use, such as Gemini, Claude and GPT, are large language models (LLMs) over billions or trillions of parameters, usually run on massive datacenters. On the other hand, there are a few small language models capable of running on laptops or mobile devices. The dissertation focuses on such a small language model, Qwen~3--0.6~B, the smallest in the Qwen~3 family of models from Alibaba \cite{yang2025qwen3technicalreport}.
  \item[Constitution.] A constitution is a written set of principles or rules defined in order to bring certain characteristics in an entity. The dissertation contains a written set of twenty-five principles (P~1 to P~25) defined in Chapter~\ref{chap:methodology}, Section~\ref{subsec:mr-constitution}, which is grounded in Anthropic's Constitutional AI framework \cite{bai2022constitutionalaiharmlessnessai}. The complete set of the constitution is added in Appendix~\ref{chap:constitutional-principles-reference}.
  \item[Constitutional distillation.] Distillation trains a small student model on data produced by a larger teacher (the supervised fine-tuning stage); constitutional distillation embeds the principles in those examples so that they are encoded in the weights rather than referenced from a document \cite{bai2022constitutionalaiharmlessnessai}.
  \item[Trustworthy.] A response is treated as trustworthy when it complies with the constitution's principles and the user can interrogate why a given answer was produced.
  \item[Personalised.] A response is personalised when information about the user, retained across conversational turns, is used to adapt it; context is gathered through the 5W+H framework with first-principles reasoning before the model answers.
  \item[Privacy-preserving.] No user data leaves the device and no cloud API is called at inference time with the user's data \cite{euhleg2019}.
  \item[Alignment tax and safety tax.] The established names for the trade-off this dissertation measures below one billion parameters: alignment or safety fine-tuning that strengthens one behaviour (constitutional compliance) does so at a measurable cost to another (visible step-by-step reasoning), so training redistributes capability rather than adding it. The general cost that alignment training imposes on capability is the \emph{alignment tax} \cite{askell2021generallanguageassistantlaboratory, ouyang2022traininglanguagemodelsfollow}; the specific case in which safety alignment degrades reasoning is the \emph{safety tax} \cite{huang2025safetytaxsafetyalignment, fu2026reducingsafetytaxllm} (Chapters~\ref{chap:methodology} and~\ref{chap:discussion}).
  \item[Thinker--Executor.] The two-model architecture of Experiment~3: where one 0.6B model (the Thinker) trained for constitutional reasoning and planning, and a second 0.6B model (the Executor) trained for tool execution, run sequentially on-device so that a single model computes at any instant. This creates an orchestration between thinking and executing the tasks for the user.
  \item[User model.] The structured, persistent record of what the system believes about the user, stored on-device as a graph of beliefs in five categories (goals, tasks, skills, styles, contexts), with a correction edge that lets the user contest a belief without the original being deleted; populated through 5W+H questioning.
  \item[5W+H.] 5W+H is a mental model to provide a framework of asking 5W (What, Where, When, Who, Why) and a How to a given situation. The dissertation uses this in both prompt instructions to the model and a bounded note which carries across conversational turns holding the 5W+H of the user's request.
\end{description}

Technical terms and state-of-the-art techniques are later defined and discussed in Chapter~\ref{chap:background} and Chapter~\ref{chap:state-of-the-art}.

% ==================================================================================================
% [DRAFT 2026-07-25 | reader retention | DEVICE B -- THE RUNNING EXAMPLE, with exact insertion points]
%
% PROBLEM: the document currently uses a different throwaway example in each chapter -- a quick
% workout, a book recommendation, Apple's market cap, "2+2 is", a job offer in Gothenburg. Each works
% locally; none accumulates. A reader crossing a hundred pages of shifting abstractions has nothing
% concrete to hold on to, and that is a large part of what "getting lost" actually is.
%
% FIX: thread ONE request through every chapter. Recommended: "Can you recommend a good programming
% book for me?" -- it is already Case~2 (Table~\ref{tab:case-p6}), it needs no technical background,
% all five conditions answer it differently, and the correct behaviour (ask one question before
% advising) is precisely the ask-before-assuming move the constitution is built around. The reader
% meets it in Chapter~1 as a puzzle and sees it resolved in Chapter~5 as data.
%
% Cost is near zero: each insertion is one or two sentences and mostly REPLACES an existing example.
% The five insertion points, in order:
%
% (1) HERE, at the end of Section~\ref{sec:definitions}, before the Research Question. INSERT:
%     "One request will recur throughout this dissertation as a worked illustration. A user asks,
%      \emph{``Can you recommend a good programming book for me?''} The question is easy to answer
%      badly and hard to answer well: a good recommendation depends entirely on who is asking, and
%      the assistant has not been told. Whether a small model recognises that, and asks before it
%      advises, turns out to separate the five systems compared in this study more sharply than any
%      aggregate score does (Table~\ref{tab:case-p6})."
%
% (2) Background, Section~\ref{subsec:why-trustworthiness}. APPEND to the paragraph beginning
%     "Returning to the recurring request, this is why a model asked to recommend a programming book
%      will happily name one: nothing in next-token prediction distinguishes a recommendation it has
%      grounds for from one it has merely made fluent."
%
% (3) Background, Section~\ref{sec:personalisation-and-user-modelling}. The live paragraph already
%     opens with the ``quick workout'' illustration, which is vivid and should stay. APPEND one
%     clause to tie it to the thread:
%     "...or maybe an old person --- and the same is true of the recurring request for a programming
%      book, where the right answer for a beginner and for a career engineer share almost nothing."
%
% (4) Methodology, Section~\ref{subsec:mr-constitution}, after the tier discussion. APPEND:
%     "The recurring book request is governed by P6, the user-context gate, and is probed directly:
%      asked with no personal context available, the trustworthy response is a single targeted
%      question, and answering anyway is scored as a failure however good the recommendation."
%
% (5) Discussion, Section~\ref{subsec:h5-thinker-executor} -- ALREADY PRESENT, no edit needed. The
%     case-study paragraph there discusses exactly this question. Once (1)--(4) are in place, add the
%     two words "the recurring" before "context-gate case" so the reader registers the loop closing.
%
% If a second thread is wanted for the honesty principles rather than the personalisation ones,
% Case~1 (the market-cap question, Table~\ref{tab:case-p5}) works the same way. Use one or two, not
% five -- the point of a running example is that it is the same one.
% ==================================================================================================
\section{Research Question}
\label{sec:research-question}

This dissertation asks the following research question: ``Can a framework teach a small language model the principles of a written constitution, and serve it through a deterministic on-device layer of prompts and tools, such that the model delivers trustworthy, personalised and privacy-preserving assistance entirely on-device?''

Note that the goal is to create a framework for developing trustworthiness in small models rather than competing against the large frontier models. For quick ideation and experimentation the dissertation focuses on the Qwen~3--0.6~B model.

\section{Hypotheses}
\label{sec:hypotheses}

This dissertation will test the below five hypotheses based on the research question formulated above; these are introduced below:

\textbf{First hypothesis} (the compliance improvement survives in a small model). Constitutional distillation produces a measurable improvement in behavioural compliance across the probed constitutional principles relative to the vanilla baseline, and the improvement holds rather than vanishing for want of capacity. This adapts the strategy of Zhang \cite{zhang2025constitutioncollapseexploringconstitutional}.

\textbf{Second hypothesis} (the improvement has a capacity cost). The same distillation that improves compliance measurably degrades reasoning capacity in a small model, observable as a collapse in \texttt{<think>}-trace length and a rise in the blank-trace rate. Under a fixed parameter budget, compliance and reasoning trade off against one another.

\textbf{Third hypothesis} (prompt engineering with tools underperforms constitutional fine-tuning). An untrained model given a carefully engineered system prompt and full tool access recovers less constitutional compliance than a model with the constitution encoded into its weights. This is the hypothesis the tool-equipped vanilla baseline exists to test.

\textbf{Fourth hypothesis} (the failure boundary is systematic, not random). Where the constitutionally fine-tuned single model underperforms the vanilla baseline, the underperformance is predictable from task type and prompt structure rather than randomly distributed, as seen by Rainone et al.\ \cite{rainone2025replacingthinkingtoolusage} and Wei et al.'s Beyond ReAct \cite{wei2025reactplannercentricframeworkcomplex} where they observe instability in responses.

\textbf{Fifth hypothesis} (separating thinking from execution restores capacity at small scale). At 0.6~B a single model must serve constitutional reasoning, tool execution and compliance from one parameter budget, with its attention over every context divided across all three demands at once. Splitting the work should relieve that contention: a Thinker model whose entire budget and attention serve constitutional reasoning alone, and an Executor model whose budget serves tool execution alone, both 0.6~B and both on-device. If the division of capacity is what displaces reasoning, then two specialised models should together recover capability that a single 0.6~B model cannot reach.

\section{Research Experiments}
\label{sec:research-experiments}

Figure~\ref{fig:architecture} gives the whole system at a glance: a small on-device model, taught a written constitution, answers a fixed set of test questions, and a larger, independent model grades the answers against a rubric; the same model is then built three ways, one per experiment. To answer the research question and test the hypotheses, the framework for this dissertation is developed by doing three step-by-step experiments, each motivated by the current research going on in the space of state of the art large language model development.

\begin{figure}[tbp]
  \centering
  \begin{tikzpicture}[>=Latex, font=\footnotesize,
      box/.style  ={draw, rounded corners, fill=black!4, align=center, inner sep=5pt},
      model/.style={draw, rounded corners, fill=tcd_blue, text=white, align=center, inner sep=6pt, minimum width=4.2cm},
      tool/.style ={draw, rounded corners, fill=tcd_blue!15, align=center, inner sep=5pt, minimum width=4.2cm},
      exp/.style  ={draw, rounded corners, fill=black!4, align=center, text width=3.3cm, minimum height=1.5cm, inner sep=5pt},
      lbl/.style  ={font=\scriptsize\itshape, align=center}]
    % --- TOP: the written constitution, encoded during training ---
    \node[box, text width=5.4cm] (con) at (5.8,7.0)
      {\textbf{Constitution}\\[2pt]\scriptsize 25 principles (P1--P25) for honest, careful, personalised answers};
    % --- CENTRE: the on-device model and its tools ---
    \node[model] (mdl) at (5.8,4.2) {\textbf{Small language model}\\ Qwen3-0.6B};
    \node[tool]  (tls) at (5.8,2.7) {Tools: web search, calculator};
    \draw[<->, thick] (mdl) -- (tls);
    \node[draw, dashed, rounded corners, thick, fit=(mdl)(tls), inner sep=10pt] (dev) {};
    \node[anchor=south west, font=\scriptsize\bfseries] at (dev.north west) {ON DEVICE};
    \draw[->, thick] (con) -- node[lbl, right, xshift=3pt]{encoded during training\\ (constitutional distillation)} (mdl);
    % --- LEFT: the fixed test questions go into the model ---
    \node[box, text width=2.4cm] (tq) at (0.2,4.2)
      {\textbf{Test questions}\\[2pt]\scriptsize the same fixed set for every condition};
    \draw[->, thick] (tq) -- (mdl);
    % --- RIGHT: an independent judge grades the answers against a rubric ---
    \node[box, text width=2.6cm] (judge)  at (11.0,4.2) {\textbf{Judge LLM}\\[2pt]\scriptsize a large, independent model};
    \node[box, text width=2.6cm] (rubric) at (11.0,2.7) {\textbf{Rubric}\\[2pt]\scriptsize what a good answer looks like};
    \node[box, text width=2.2cm] (grade)  at (11.0,6.0) {\textbf{Grade}};
    \draw[->, thick] (mdl) -- node[lbl, above]{answers} (judge);
    \draw[->, thick] (rubric) -- (judge);
    \draw[->, thick] (judge) -- (grade);
    % --- BOTTOM: the same small model was built three ways, one per experiment ---
    \node[draw, rounded corners, fill=tcd_blue!10, font=\bfseries, align=center,
          minimum width=8cm, minimum height=0.6cm] (hd) at (5.8,1.15) {The small model was built three ways};
    \node[exp] (exp1) at (1.9,-1.0) {\textbf{Experiment 1}\\[2pt]\scriptsize Taught the format of a good answer (template SFT)};
    \node[exp] (exp2) at (5.8,-1.0) {\textbf{Experiment 2}\\[2pt]\scriptsize Taught the behaviour of a large teacher (constitutional distillation)};
    \node[exp] (exp3) at (9.7,-1.0) {\textbf{Experiment 3}\\[2pt]\scriptsize Split in two: one thinks, one acts (Thinker--Executor)};
    \draw[thick] (dev.south) -- (hd.north);
    \draw[thick] (hd.south) -- (5.8,0.4);
    \draw[thick] (1.9,0.4) -- (9.7,0.4);
    \draw[->, thick] (1.9,0.4) -- (exp1.north);
    \draw[->, thick] (5.8,0.4) -- (exp2.north);
    \draw[->, thick] (9.7,0.4) -- (exp3.north);
  \end{tikzpicture}
  \caption{The measurement apparatus and the three experiments. A small on-device model (Qwen3-0.6B), taught a written twenty-five-principle constitution, answers a fixed set of test questions; a larger, independent model then judges the answers against a rubric and assigns a grade. The same small model is built three ways, one per experiment: template supervised fine-tuning, constitutional distillation from a frontier-scale teacher, and a Thinker--Executor split into two on-device 0.6B models.}
  \label{fig:architecture}
\end{figure}

\textbf{Experiment 1} applies template-based supervised fine-tuning to constitutionally annotated examples, using a 16-bit LoRA \cite{hu2021loralowrankadaptationlarge} on Qwen~3--0.6~B. It replicates the SL-CAI stage of Bai et al.\ \cite{bai2022constitutionalaiharmlessnessai} and the constitutional SFT strategy of Zhang \cite{zhang2025constitutioncollapseexploringconstitutional} at roughly one-thirteenth the scale. This is the simplest possible test of whether constitutional training produces a measurable compliance change in a small model.

\textbf{Experiment 2} changes exactly one component over Experiment 1: it replaces the template-based dataset with one constructed under the Anthropic Constitutional AI framework \cite{bai2022constitutionalaiharmlessnessai}, distilled from a frontier-scale teacher model, over the same base weights and the same training recipe. This tests whether Experiment~1's shortfall is a data-quality problem or a scale problem, and, read against the tool-equipped baseline, it carries the third hypothesis's comparison between prompt engineering and constitutional fine-tuning.

\textbf{Experiment 3} replaces the single generalist model with a Thinker--Executor dual-SFT architecture: one 0.6~B model trained for constitutional reasoning, one for tool execution. The task is decomposed so that each model owns one capability, devoting its full parameter budget and its attention over a single-role context to that capability alone. This extends the current research on the planner--executor pattern of Beyond ReAct \cite{wei2025reactplannercentricframeworkcomplex} and Molinari and Ciravegna's Reason-Plan-ReAct \cite{molinari2025reasonplanreactreasonerplannersupervisingreact}, with the difference that the large cloud executor is now replaced by a second on-device 0.6~B model.

Chapter~\ref{chap:methodology} specifies the architecture for all three experiments and Chapter~\ref{chap:experiments-results} reports the results of each in turn.

\section{Scope and Assumptions}
\label{sec:scope-and-assumptions}

This dissertation evaluates only one model (Qwen~3--0.6~B), one training method (16-bit LoRA supervised fine-tuning), a synthetic dataset and one benchmark framework (the constitutional principles P~1 to P~25). Qwen~3--0.6~B was chosen for three main reasons: it supports a native \texttt{<think>} block in the base model, it can perform tool calls, and it is compatible with Unsloth for single-consumer-GPU training \cite{yang2025qwen3technicalreport}.

Moreover, 16-bit LoRA is chosen over full fine-tuning because it adapts the model without overwriting pre-trained representations, which Hu et al.\ established as the standard practice in parameter-efficient alignment at this scale \cite{hu2021loralowrankadaptationlarge}. 

% [APPLIED 2026-07-25] The Rainone misattribution was corrected in the live text at all six
% sites (introduction x2, background, state-of-the-art, conclusion, discussion x2). RAGEN
% \cite{wang2025ragenunderstandingselfevolutionllm} now carries the 0.5B RL-collapse evidence and
% Rainone is cited only for what it shows. The instruction block below is retained as the record of
% what was wrong and why; do NOT re-apply it.
% [AUDIT-FIX 2026-07-24 | DROP-IN] ACTION: REPLACE the live sentence flagged in the COPY-PASTE note in this block (the "Rainone et al. find GRPO unstable at 0.6B" attribution), per that note; apply the SAME swap at Conclusion Section~\ref{subsec:conclusion-future-grpo} and Background Section~\ref{sec:supervised-fine-tuning-and-parameter-efficient-adaptation}.
% [AUDIT-FIX 2026-07-24 -- citation misattribution; the SAME swap applies at Conclusion
% Section~\ref{subsec:conclusion-future-grpo} and Background Section~\ref{sec:supervised-fine-tuning-and-parameter-efficient-adaptation}.
% "Rainone et al. find GRPO unstable at 0.6B" is wrong: their smallest model is Llama-1B, and their result
% is the OPPOSITE -- structured tool-use (Chain-of-Edits) makes RL WORK at 1B-3B. The 0.5B RL-collapse
% evidence is RAGEN (Echo Trap on Qwen2.5-0.5B), ALREADY in ref.bib as
% \cite{wang2025ragenunderstandingselfevolutionllm} (currently uncited), plus Beyond ReAct (already cited).
% COPY-PASTE: replace the sentence below --
%   "...excluded because Rainone et al.\ find GRPO empirically unstable at 0.6~B and attribute this to model
%    scale \cite{rainone2025replacingthinkingtoolusage}, and Beyond ReAct explicitly excludes 0.6~B from its
%    RL phase for the same reason \cite{wei2025reactplannercentricframeworkcomplex}"
% -- with --
%   "...excluded because reinforcement learning is empirically unstable below one billion parameters: RAGEN
%    reports multi-turn RL collapse on a 0.5B model \cite{wang2025ragenunderstandingselfevolutionllm}, and
%    Beyond ReAct excludes Qwen3-0.6B from its RL phase for the same reason
%    \cite{wei2025reactplannercentricframeworkcomplex}"
% Keep \cite{rainone2025replacingthinkingtoolusage} only for the distinct point it actually supports
% (small-model reasoning brittleness, recoverable via structured tool use -- e.g. Hypothesis 4,
% Section~\ref{sec:hypotheses}), and drop the "attributed to scale" phrasing throughout.]
There are three areas which are excluded from this dissertation, due to resource and time limitations. GRPO and reinforcement learning post-training are excluded because reinforcement learning is empirically unstable below one billion parameters: RAGEN reports multi-turn reinforcement-learning collapse on a 0.5~B model \cite{wang2025ragenunderstandingselfevolutionllm}, and Beyond ReAct explicitly excludes Qwen~3--0.6~B from its reinforcement-learning phase for the same reason \cite{wei2025reactplannercentricframeworkcomplex}. In addition, GRPO takes a huge amount of compute capacity and time, which was not possible. Moreover, the automated probes may not capture what real users perceive as trustworthy, especially in the multi-turn personalisation which will need a large number of users to test the model efficiency and will need a separate pipeline for an ethical framework to delete the user data, so a human-evaluation protocol will be added as the primary future-work direction. Models above 1~B are excluded, but not because they breach the privacy guarantee, which on-device execution preserves at any model size. They are excluded because the device tier this dissertation targets, the low-end and older handsets that flagship-only deployments such as Apple's AFM and Google's Gemma~4 leave behind, cannot run a model above 1~B within a sustained thermal and energy budget once the reasoning traces and tool calls are accounted for, as Tummalapalli et al.\ show \cite{tummalapalli2026llminferenceedgemobile}.

\section{Contributions}
\label{sec:contributions}

This dissertation's primary contribution is a model-agnostic method for measuring and improving constitutional behaviour at small-model scale, rather than a model that competes with a large one or a simple yes-or-no verdict on whether small models can ``comply'' with a constitution. The method currently validates on the single Qwen~3--0.6~B model, owing to hardware constraints and the need for rapid experimentation, with cross-model validation left for future work (Chapter~\ref{chap:conclusion}). Its second contribution is architectural: a Thinker--Executor design in which both components are on-device small models, replacing the large cloud executor of small-agent collaboration designs such as Zywot et al.'s \cite{zywot2026smallagentscollaboratebeat} with a second 0.6~B model, so that the whole system can run on the user's own device. The contributions are set out in detail in Section~\ref{sec:conclusion-contributions}.

% ==================================================================================================
% [DRAFT 2026-07-25 | reader comprehension -- the single highest-value addition in the document]
% THE ARGUMENT OF THE DISSERTATION, ON ONE PAGE.
%
% WHY THIS IS NEEDED. The dissertation contains a genuinely tight causal argument -- privacy forces
% the model onto the device, the device forces it below one billion parameters, that size breaks the
% self-critique step canonical Constitutional AI depends on, which forces the asymmetric teacher, which
% buys compliance at the cost of reasoning, which reveals a fixed capacity budget, which motivates
% splitting the budget across two models. Every link in that chain is established somewhere in the
% document. NOWHERE IS THE CHAIN WRITTEN DOWN AS A CHAIN. A reader has to reconstruct it across seven
% chapters and a hundred pages, and most readers will not.
%
% Section~\ref{sec:dissertation-structure} below tells the reader what each chapter CONTAINS. It does
% not tell them what the dissertation ARGUES. Those are different jobs and the second one is missing.
% Placing the chain here means every later chapter is read as a step the reader is already expecting,
% which is what "connecting the dots" actually requires: the dots have to be visible before the reader
% meets them, not only after.
%
% Place immediately before \section{Dissertation Structure}, under a heading such as
% \section{The Argument in Brief}. Keep it as a numbered chain -- the numbering is what makes the
% dependency visible; as flowing prose it degrades back into a summary. VERIFY each step is one you
% are willing to defend, and adjust step 10 if the results are re-read.
%
% \section{The Argument in Brief}
% \label{sec:argument-chain}
%
% The chapters that follow develop a single connected argument, and it is set out here in full so that
% each chapter can be read as the step it is rather than as a self-contained topic.
%
% \begin{enumerate}
%   \item An assistant asked about health, money and personal circumstances accumulates a record not
%     of what a person did but of how they reasoned, and that record is the special-category data
%     Article~9 of the General Data Protection Regulation singles out for protection
%     (Section~\ref{sec:motivation}).
%   \item The usual safeguards govern what a provider stores, not what a model has already absorbed
%     and can be induced to repeat; federated learning removes the transmission but leaves the residue
%     in the weights (Section~\ref{sec:motivation}).
%   \item The data must therefore stay on the user's own device, which is a constraint on the
%     architecture rather than on the contract (Section~\ref{subsec:gdpr-and-the-case-for-local-inference}).
%   \item The device tier this work targets --- the low-end and older handsets that flagship on-device
%     deployments leave unserved --- bounds the model below one billion parameters
%     (Section~\ref{sec:scope-and-assumptions}).
%   \item At that size the model cannot reliably critique its own answers, and self-critique is the
%     step on which canonical Constitutional AI depends; it already degrades at eight billion
%     parameters (Section~\ref{sec:constitutional-ai}, Section~\ref{sec:sota-cai-scale}).
%   \item The constitution must therefore be taught rather than consulted: a frontier teacher reasons
%     with the principles in view, and the student learns only from the behaviour that results, never
%     seeing the rules as text (Section~\ref{subsec:mr-dataset}).
%   \item Taught this way, compliance rises to the highest of the five conditions measured --- and the
%     model's visible reasoning collapses (Section~\ref{subsec:mr-exp2-results}).
%   \item That pairing is not a defect of the training run but the small-model form of the established
%     \emph{alignment tax}: at this scale capacity behaves as a fixed budget that training
%     redistributes rather than enlarges (Section~\ref{sec:capacity-displacement-trade-off}).
%   \item If the budget is fixed, the remaining structural move is to divide it, so one model reasons
%     and another acts, both on the device and run in sequence so that only one computes at any instant
%     (Section~\ref{sec:mr-exp3}).
%   \item The division restores the visible reasoning and produces the only condition that asks before
%     assuming, but not the headline compliance and not long-conversation stability: trustworthiness at
%     this scale is bought in layers, and no single layer buys all of it
%     (Section~\ref{subsec:h5-thinker-executor}, Section~\ref{sec:implications-practice}).
%   \item What therefore survives the next release of a better small model is not the verdict on this
%     one but the instrument that located the boundary, which is why the framework rather than the
%     result is offered as the contribution (Section~\ref{subsec:transferable-framework}).
% \end{enumerate}
%
% OPTIONAL, and strong if the space allows: render steps 1--10 as a single vertical flow diagram, each
% box one step and each arrow labelled with the word that forces the next ("therefore", "but", "so").
% A reader who sees the argument as a picture before reading it as prose retains it far better, and it
% would give the Introduction the figure it currently lacks. If only one figure goes into this chapter,
% though, prefer fig:architecture --- it is already drafted and seven other openers reference it.
% ==================================================================================================
\section{Dissertation Structure}
\label{sec:dissertation-structure}

In order to build the foundation, and to test the hypotheses through experiments, this dissertation is organised as follows.

\begin{enumerate}
    \item Chapter~\ref{chap:background} introduces various background and literature to build the foundation for reviewing key concepts around the very basics of language models, constitutional AI, harness and scaffolding approaches, and analysing the gap for small-model brittleness and how a Thinker--Executor Framework might help in achieving delegations.
    \item Chapter~\ref{chap:state-of-the-art} surveys where the field stands today, the on-device flagship models, constitutional alignment as it has been shown at larger scales, planner-executor decomposition, and personalisation memory, and closes on the gap bringing a small on-device system that is both constitutionally trained and tool-grounded.
    \item Chapter~\ref{chap:methodology} presents the methodology for experimenting the three-experiment pipeline as a progressive design; this chapter shall build on various techniques designed for the model training architecture to perform the experiments. 
    \item Chapter~\ref{chap:experiments-results} then reports what each of the three experiments found, showcasing the contribution of each experiment.
    \item Chapter~\ref{chap:discussion} interprets those results against the five hypotheses, compares them with literature from Chapter~\ref{chap:background}, along with the limitations from Section~\ref{sec:scope-and-assumptions} and answers the research question, showing where the methodology passes and where this approach is limited.
    \item Chapter~\ref{chap:conclusion} answers the research question, concluding on what is possible with this framework to present a Trustworthy Personalised AI, along with future work in the ever-growing field of language models.
\end{enumerate}

% [DRAFT 2026-07-25 | reader retention | DEVICE A -- CHAPTER-CLOSING BRIDGE]
% Three sentences: what to carry forward, the question now open, the chapter that answers it. An
% open question is the actual reason a reader turns the page; a chapter that simply stops does not
% give them one. Append as the last paragraph of the chapter. To apply, delete the "% " prefix.
%
% Everything set out above is a claim about what ought to be possible. Whether any of it is possible depends on how a language model actually works, and in particular on what a model small enough to run on a telephone can and cannot hold at once. That is the work of the next chapter, which establishes not a survey for its own sake but the specific properties of these models that force every design decision taken later in this dissertation.
